% =====================================================================
%  باب الاستعاذة — الأبيات ٩٥–٩٩
% =====================================================================
\chapter{باب الاستعاذة}

\textbf{الاستعاذة}: طلب العَوذ، أي الالتجاء إلى الله والاحتماء به. والمراد هنا قول القارئ «أعوذ بالله من الشيطان الرجيم» قبل القراءة.
والاستعاذة \textbf{ليست من القرآن} بإجماع العلماء، ولذلك ليس فيها رموز لقراء، بل حكمها عام لجميع القراء. وعدد أبيات الباب خمسة.

\bayt{95}{إِذَا مَا أَرَدْتَ الدَّهْرَ تَقْرَأُ فَاسْتَعِذْ}{جِهَارًا مِنَ الشَّيْطَانِ بِاللهِ مُسْجَلَا}
\bayt{96}{عَلَى مَا أَتَى فِي النَّحْلِ يُسْرًا وَإِنْ تَزِدْ}{لِرَبِّكَ تَنْزِيهًا فَلَسْتَ مُجَهَّلَا}
\begin{mufradat}
\kalima{أردت}{قصدت.}
\kalima{الدهر}{في أي وقت من الزمان.}
\kalima{جِهارًا}{جهرًا، أي بصوت مسموع.}
\kalima{مُسْجَلًا}{مطلقًا: لكل القراء وفي كل القرآن.}
\kalima{على ما أتى في النحل}{بلفظ آية النحل: \ayah{فَإِذَا قَرَأۡتَ ٱلۡقُرۡءَانَ فَٱسۡتَعِذۡ بِٱللَّهِ مِنَ ٱلشَّيۡطَٰنِ ٱلرَّجِيمِ} \src{النحل: ٩٨}.}
\kalima{يُسرًا}{حال كونه سهلًا قليل الكلمات.}
\kalima{تنزيهًا}{تقديسًا وتعظيمًا لله.}
\kalima{مجهَّلًا}{منسوبًا إلى الجهل.}
\end{mufradat}

\maana[(البيتان ٩٥–٩٦)]
إذا أردت أن تقرأ القرآن في أي وقت، ولأي قارئ، ومن أي موضع (أول سورة أو وسطها)، فاستعذ \textbf{جهرًا} قبل القراءة، بلفظ آية النحل:
\begin{center}
\textbf{\large أَعُوذُ بِاللهِ مِنَ الشَّيْطَانِ الرَّجِيمِ}
\end{center}
وهذا هو اللفظ المختار لسهولته وموافقته للقرآن. وإن زدت عليه ألفاظًا فيها تعظيم لله فلا حرج، ولا تُنسب إلى الجهل.

\begin{amthila}[صيغ جائزة فيها زيادة تنزيه]
\begin{itemize}
\item أعوذ بالله \textbf{العظيم} من الشيطان الرجيم.
\item أعوذ بالله \textbf{السميع العليم} من الشيطان الرجيم.
\end{itemize}
\end{amthila}

\begin{qaida}[لماذا قال «إذا ما أردت»؟]
الآية تقول: \ayah{فَإِذَا قَرَأۡتَ ٱلۡقُرۡءَانَ فَٱسۡتَعِذۡ}، وظاهرها أن الاستعاذة بعد القراءة. فبيّن الناظم أن معناها: \textbf{إذا أردتَ} القراءة فاستعذ، أي قبلها. ونظيره: \ayah{إِذَا قُمۡتُمۡ إِلَى ٱلصَّلَوٰةِ فَٱغۡسِلُواْ} \src{المائدة: ٦} أي إذا أردتم القيام إلى الصلاة.
\end{qaida}

\bayt{97}{وَقَدْ ذَكَرُوا لَفْظَ الرَّسُولِ فَلَمْ يَزِدْ}{وَلَوْ صَحَّ هَذَا النَّقْلُ لَمْ يُبْقِ مُجْمَلَا}
\begin{mufradat}
\kalima{ذكروا}{أي بعض العلماء من القراء والمحدثين.}
\kalima{لفظ الرسول}{لفظ استعاذة النبي ﷺ.}
\kalima{مُجمَلًا}{إجمالًا وإطلاقًا في الآية.}
\end{mufradat}
\maana ذكر بعض العلماء أحاديث فيها أن النبي ﷺ لم يزد على لفظ آية النحل في الاستعاذة. ولو صحت هذه الأحاديث لكان اللفظ متعينًا ولم يبق في الآية إجمال، فلا تجوز الزيادة. لكنها أحاديث \textbf{ضعيفة}، وجاءت أحاديث أصح منها فيها زيادة، مثل: «أعوذ بالله السميع العليم من الشيطان الرجيم من همزه ونفخه ونفثه» (رواه أبو داود والترمذي). فبقيت الآية على إطلاقها، فتجوز الزيادة والنقص.

\begin{tanbih}
الأمر بالاستعاذة في الآية \textbf{للندب} (الاستحباب) عند جمهور العلماء من السلف والخلف، لا للوجوب.
\end{tanbih}

\bayt{98}{وَفِيهِ مَقَالٌ فِي الْأُصُولِ فُرُوعُهُ}{فَلَا تَعْدُ مِنْهَا بَاسِقًا وَمُظَلِّلَا}
\begin{mufradat}
\kalima{مقال}{كلام كثير.}
\kalima{الأصول}{كتب أصول الفقه والحديث والقراءات.}
\kalima{فروعه}{تفاصيله (شبّهها بأغصان الشجرة).}
\kalima{لا تَعْدُ}{لا تتجاوز.}
\kalima{باسقًا ومظلِّلًا}{الشجر الطويل الوارف الظل، والمراد: القول القوي المؤيَّد بالدليل.}
\end{mufradat}
\maana في التعوذ كلام طويل في كتب الأصول: هل الأمر للوجوب أو الندب؟ وما درجة الأحاديث فيه؟ وما حكم الجهر والإخفاء والوقف والوصل؟ فارجع إليها ولا تتجاوز القول القوي الذي تؤيده الأدلة.

\bayt{99}{وَإِخْفَاؤُهُ فَصْلٌ أَبَاهُ وُعَاتُنَا}{وَكَمْ مِنْ فَتًى كَالْمَهْدَوِي فِيهِ أَعْمَلَا}
\begin{mufradat}
\kalima{إخفاؤه}{الإسرار بالتعوذ.}
\kalima{فصل}{فرق (بين القرآن وغيره).}
\kalima{أباه}{رفضه.}
\kalima{وُعاتنا}{علماؤنا الحفاظ (جمع واعٍ).}
\kalima{المهدوي}{الإمام أبو العباس أحمد بن عمار المهدوي المقرئ المفسر (ت بعد ٤٣٠ هـ).}
\kalima{أعمل}{أعمل فكره في تصحيحه.}
\end{mufradat}
\maana إخفاء التعوذ (عند من قال به) فيه حكمة: الفرق بين ما هو قرآن وما ليس بقرآن. لكن علماءنا الحفاظ رفضوا الإخفاء وأخذوا بالجهر في جميع القرآن لكل القراء. وكم من عالم كالمهدوي أعمل فكره في تصحيح الإخفاء والإقراء به.

\begin{tanbih}[هل في البيت رمز؟]
حمل بعض الشراح الفاء في «فصل» على أنها رمز \Hamza، والألف في «أباه» على أنها رمز \Nafi، أي أنهما يخفيان التعوذ. والصحيح (كما في الوافي) أنه \textbf{لا رمز في البيت}؛ لأن الباب ليس فيه خلاف بين القراء، والعمل على \textbf{الجهر للجميع}. وقد رُويت عن \Hamza{} و\Nafi{} روايات بالإخفاء من غير طريق الشاطبية.
\end{tanbih}

\begin{qaida}[متى يُستحب إخفاء التعوذ؟ (من الوافي)]
يُستحب الجهر بالتعوذ في عموم الأحوال، ويُستحب الإخفاء في أربعة مواضع:
\begin{enumerate}
\item إذا كان القارئ يقرأ سرًّا، منفردًا أو في مجلس.
\item إذا كان خاليًا (وحده) سواء قرأ سرًّا أو جهرًا.
\item إذا كان في الصلاة (سرية أو جهرية، إمامًا أو مأمومًا أو منفردًا).
\item إذا كان يقرأ وسط جماعة يتدارسون القرآن ولم يكن هو المبتدئ بالقراءة.
\end{enumerate}
وفائدة الجهر: أن يُنصت السامع للقراءة من أولها فلا يفوته منها شيء.
\end{qaida}

\begin{tahrir}
\ithaf{6}{إِذَا مَا أَرَدْتَ الدَّهْرَ تَقْرَأُ فَاسْتَعِذْ}{وَبِالْجَهْرِ عِنْدَ الْكُلِّ فِي الْكُلِّ مُسْجَلَا}
\ithaf{7}{بِشَرْطِ اجْتِمَاعٍ وَابْتِدَاءِ دِرَاسَةٍ}{وَلَا مُخْفِيًا أَوْ فِي الصَّلَاةِ فَفَصِّلَا}
\ithaf{8}{وَوَقْفٌ عَلَيْهِ ثُمَّ وَصْلٌ بِأَرْبَعٍ}{لَهُمْ وَاسْتَعِذْ نَدْبًا أَوَ اوْجِبْ وَوَهِّلَا}
\medskip
\textbf{الشرح:}
\begin{itemize}
\item الجهر بالتعوذ لكل القراء في كل القرآن، \textbf{بشرط}: أن يكون معه من يسمعه، وأن يكون هو المبتدئ بالقراءة في مجلس الدرس، وألا يكون يقرأ سرًّا، وألا يكون في الصلاة. فإن اختلّ شرط أخفاه. (وهذا موافق للمواضع الأربعة المذكورة في الوافي.)
\item يجوز \textbf{الوقف} على التعوذ، ويجوز \textbf{وصله} بما بعده، لكل القراء.
\item \textbf{بأربع}: إذا اجتمعت الاستعاذة والبسملة وأول السورة ففيها \textbf{أربعة أوجه} لكل القراء (انظرها في الجدول أدناه).
\item الاستعاذة \textbf{مندوبة}، وقيل واجبة، وهذا القول الثاني ضعيف (معنى «وهِّلا»: اعتبره وهمًا).
\end{itemize}
\end{tahrir}

\begin{khulasa}[الأوجه الأربعة عند اجتماع الاستعاذة والبسملة وأول السورة]
\centering
\renewcommand{\arraystretch}{1.5}
\begin{tabular}{c l l}
\toprule
\textbf{الوجه} & \textbf{الاسم} & \textbf{الصورة}\\
\midrule
١ & قطع الجميع & أعوذ بالله... \textbf{(وقف)} بسم الله... \textbf{(وقف)} الحمد لله...\\
٢ & قطع الأول ووصل الثاني بالثالث & أعوذ بالله... \textbf{(وقف)} بسم الله... الحمد لله...\\
٣ & وصل الأول بالثاني وقطع الثالث & أعوذ بالله... بسم الله... \textbf{(وقف)} الحمد لله...\\
٤ & وصل الجميع & أعوذ بالله... بسم الله... الحمد لله...\\
\bottomrule
\end{tabular}
\end{khulasa}

\begin{tanbih}[تتمة: إذا قطع القارئ قراءته]
لو قطع القارئ قراءته لأمر طارئ قهري (كعطاس أو تنحنح) أو لكلام يتعلق بالقراءة (كسؤال عن شيء فيها)، فإنه \textbf{لا يعيد} التعوذ.
أما لو قطعها إعراضًا عنها، أو لكلام لا يتعلق بها ولو ردًّا لسلام، فإنه \textbf{يستأنف} التعوذ.
\end{tanbih}
